% Cover letter using letter.cls
\documentclass{letter}
\topmargin=-1in
\textheight=8in
\oddsidemargin=0pt
\textwidth=6.5in

\begin{document}

\signature{N. Ohayon Rozanes}
\longindentation=0pt
\let\raggedleft\raggedright

\begin{letter}{}

  \begin{center}
    {\large\bf N. Ohayon Rozanes} \\
    {611 Granbury Dr. \\ Allen, Texas, 75013 \\ 469-560-0141}
  \end{center} \vfill

  \opening{Dear Hiring Manager:}

  \noindent I am writing to apply for the Software Engineering Internship at
  Twilio for Summer 2026. I am pursuing a double degree in Data Science and
  Mathematics at the University of Texas at Dallas. Most of my applied work
  has been at the boundary between systems software and real-time
  constraints --- first in embedded firmware at Exolambda, and more recently
  in data pipelines and performance-critical research code --- and I have
  been looking to take that experience into distributed systems. Twilio's
  core engineering problems sit right at that intersection.

  \noindent The bulk of my relevant experience comes from two directions. At
  Exolambda I built an embedded communication stack and sensor pipeline that
  reduced end-to-end latency by 40\%, wrote real-time signal-processing and
  control code that had to meet strict latency targets on every run, and
  generally got comfortable debugging failures at the hardware-software
  boundary under deadline pressure. That work gave me a concrete sense of
  what real-time means in practice: not just fast, but deterministically
  correct within a time budget, where a missed deadline is a functional
  failure. On the software architecture side, I built a physics simulator
  with a gRPC service layer so that Python clients could drive and script
  simulation batches across a distributed boundary, which gave me hands-on
  experience designing interfaces where correctness and latency are both
  constraints. I have also spent considerable time squeezing performance out
  of data pipelines in my bioinformatics research role, where GPU
  parallelism cut a 12-hour computation down to 3 hours, and I have come to
  think of performance work as a form of careful reasoning rather than just
  optimization.

  \noindent What makes Twilio's engineering problems interesting to me is
  the scale at which real-time constraints have to hold. In embedded work,
  real-time is hard because the hardware is unforgiving. In a virtualized
  cloud environment at Twilio's scale, real-time is hard for a completely
  different set of reasons: distributed coordination, jitter from
  virtualization, fault tolerance across components that can fail
  independently. Real-time audio at that scale is a legitimately difficult
  problem, and the failure modes are immediately visible to end users. I
  find that kind of accountability clarifying --- it is the same reason I
  liked working in aerospace, where the software either works or it does not.
  I would rather work on a constrained, high-stakes systems problem than a
  more comfortable one with softer requirements.

  \noindent I would welcome the opportunity to work on these problems with
  the Twilio team.

  \closing{Sincerely yours,}

\end{letter}

\end{document}
